% AUTORIDAD DE FUENTE: Exposé IX, §7.2 y Lema 7.2.1,
% líneas 1889--1938; índices de escaneo 413--417, folios 402--406.
% Las cinco lecturas materiales D/C se conservan mediante \sisi. También
% se conserva reversiblemente la base de étalidad S/S_n, inequívocamente
% determinada por el contexto; la cita en blanco del impreso no se inventa.

\medskip
\noindent
7.2. Consignemos ahora el siguiente lema.

\medskip
\noindent
\underline{Lema 7.2.1.} Sea $S$ el espectro de un anillo noetheriano
$V$, separado y completo para la topología definida por un ideal
$\underline m$ de $V$. Sea $C$ la categoría de los esquemas en grupos
$\Lambda$ sobre $S$, lisos y de tipo finito sobre $S$, tales que
$\Lambda/\Lambda^o$ sea un esquema en grupos finito étale sobre $S$ y
$\Lambda^o$ sea una extensión de un esquema abeliano $B$ por un toro
isotrivial $T$ (SGA~3 IX~1.1).

%%%%%%%%%%  scan idx 414  |  folio 403  %%%%%%%%%%

Sea $C'$ la categoría de los objetos en grupos de la categoría de los
esquemas formales de tipo finito (EGA~I~\sisi{10.3.3}{10.13.3}) sobre
$\operatorname{Specf}(\sisi{A}{V})$, y consideremos el funtor
\[
 A\longmapsto\widehat A
\]
de $C$ en $C'$ que toma el completado formal a lo largo de
$A_o=A\times_S S_o$, donde
$S_o=\operatorname{Spec}(V/\underline m)$. Entonces:

\smallskip
\noindent
\underline{a) El funtor precedente es plenamente fiel.}

\smallskip
\noindent
\underline{b) Sea $\mathcal A$ un objeto de $C'$, liso sobre
$\widehat S=\operatorname{Specf}(\sisi{A}{V})$}; es decir, supongamos que
los esquemas ordinarios
\[
 A_n=\mathcal A\times_{\widehat S}S_n,
 \qquad
 S_n=\operatorname{Spec}(\sisi{A}{V}/\underline m^{n+1}),
\]
son lisos sobre $S_n$. Supongamos que $A_o/A_o^o$ es un esquema en
grupos finito étale sobre $S_o$, y que $A_o^o$ es una extensión de un
esquema abeliano $B_o$ por un toro isotrivial $T_o$. Pongamos
\[
 \mathcal A^o=(A_n^o)_{n\geq0},
 \qquad
 \mathcal A/\mathcal A^o=(A_n/A_n^o)_{n\geq0}.
\]
\underline{Entonces $\mathcal A/\mathcal A^o$ es finito étale sobre
$\widehat S$, y $\mathcal A^o$ es, de manera esencialmente única y
funtorial en $\mathcal A$, una extensión de un esquema abeliano formal
$\mathcal B$ sobre $\widehat S$ por un toro formal $\mathcal T$ sobre
$\widehat S$, que se reducen respectivamente a $B_o$ y $T_o$. El esquema
formal en grupos $\mathcal A$ pertenece a la imagen esencial de
$C\longrightarrow C'$---o, como diremos también, es
``algebraizable''---si y solo si $\mathcal B$ es algebraizable.}

\medskip
\noindent
7.2.1.1. Observemos primero, sobre un esquema arbitrario $S$, que si
$A$ es un esquema en grupos liso tal que $A^o$ es una extensión de un
esquema abeliano por un toro, esta estructura de extensión es
esencialmente única y es funtorial en $A$ en el sentido evidente. En
efecto, la afirmación se reduce inmediatamente al hecho de que todo
morfismo de un toro $T$ a un esquema abeliano $B$ es nulo. Esto es bien
conocido sobre las fibras geométricas [15], y SGA~3 IX~5.3 da entonces
la conclusión.

%%%%%%%%%%  scan idx 415  |  folio 404  %%%%%%%%%%

Aplicando el resultado a los sistemas compatibles de esquemas en grupos
lisos $A_n$ sobre los $S_n$ que aparecen en 7.2.1, se obtiene el
resultado análogo para los esquemas formales en grupos sobre $S$.
Además, ya observamos en 7.1 que, si $A_o^o$ es una extensión de un
esquema abeliano por un toro, lo mismo vale para $A_n^o$. Esto demuestra
la primera afirmación de b).

\medskip
\noindent
7.2.1.2. Supongamos ahora que el esquema formal en grupos $\mathcal A$
sobre $S$ es de la forma $\widehat A$, con
$A\in\operatorname{ob}C$. Si $A^o$ es una extensión de un esquema
abeliano $B$ por un toro $T$, se tiene evidentemente
$\mathcal B\simeq\widehat B$. Así, la algebraizabilidad de $\mathcal A$
implica la de $\mathcal B$.

Recíprocamente, supongamos que $\mathcal B$ es algebraizable, es decir,
que proviene de un esquema abeliano $B$, y demostremos que $\mathcal A$
es algebraizable. Demostremos primero que
$\mathcal A^o=(A_n^o)$ es algebraizable. Por VIII~3.7, las extensiones
$A_n^o$ de $B_n$ por $T_n$ pueden interpretarse como homomorfismos
\begin{equation}\tag{*}
 \underline M_n\longrightarrow B'_n,
\end{equation}
donde $\underline M_n$ es el grupo constante torcido sobre $S_n$ dual
de $T_n$ (SGA~3 X~5.7), y $B'_n$ es el esquema abeliano dual de $B_n$.
La restricción a $S_o$, o a cualquier $S_n$, induce una equivalencia de
la categoría de los toros isotriviales sobre $S$ con la categoría
correspondiente sobre $S_o$ (SGA~3 X~3.2). Por tanto, los $T_n$
provienen de un toro $T$ sobre $S$, determinado salvo isomorfismo único.
Denotemos por $\underline M$ su grupo dual; su restricción a $S_n$ es
canónicamente isomorfa a $\underline M_n$.

La isotrivialidad de $T$ significa que $\underline M$ es una unión
disjunta de $S$-esquemas $\underline M^i$ finitos sobre $S$
(SGA~3 X~5.11). Se sigue inmediatamente que los homomorfismos $(*)$
provienen de un único homomorfismo
\begin{equation}\tag{**}
 \underline M\longrightarrow B',
\end{equation}

%%%%%%%%%%  scan idx 416  |  folio 405  %%%%%%%%%%

donde $B'$ es el esquema abeliano dual de $B$. A su vez, este
homomorfismo se interpreta como una extensión de $B$ por $T$, digamos
$A^o$, que proporciona evidentemente una algebraización de
$\mathcal A^o$.

\medskip
\noindent
7.2.1.3. Por otra parte, por hipótesis $A_o/A_o^o$ es un esquema en
grupos finito $\Phi_o$ sobre $\sisi{S}{S_o}$, y los $A_n/A_n^o$ son
asimismo esquemas en grupos finitos $\Phi_n$ sobre $\sisi{S}{S_n}$;
son necesariamente étale sobre $\sisi{S}{S_n}$. Existe, por tanto, un
esquema en grupos étale $\Phi$ sobre $S$, determinado salvo isomorfismo
único, cuyo completado formal es el esquema formal en grupos
$\mathcal A/\mathcal A^o$ definido por los $\Phi_n$.

Como $G=A^o$ es divisible, la categoría de las extensiones, como haces
fppf, de un grupo finito localmente constante sobre $S$, tal como
$\Phi$, por el esquema en grupos $A^o$ es canónicamente equivalente a
la categoría de los torsores bajo el haz
\[
 H=\underline{\operatorname{Hom}}_{\mathrm{gr}}(\Phi,G),
\]
y esta equivalencia es compatible con todo cambio de base. Aplicando la
misma descripción a las extensiones de $\Phi_n$ por $G_n$ en términos
de torsores bajo
\[
 H_n\simeq
 \underline{\operatorname{Hom}}_{\mathrm{gr}}(\Phi_n,G_n),
\]
los $A_n$ definen un sistema coherente de torsores $P_n$ bajo los
$H_n$. Como $G$ es una extensión de un esquema abeliano por un toro, se
deduce fácilmente que $H$ es un esquema en grupos \underline{finito y
localmente libre sobre $S$}. Por consiguiente, el funtor
\[
 P\longmapsto\widehat P=(P_n)
\]
es una equivalencia entre la categoría de los torsores sobre $S$ bajo
$H$ y la categoría de los torsores ``formales'' $\widehat P$ bajo
$\widehat H$, es decir, de los sistemas coherentes de torsores $P_n$
bajo los $H_n$.

%%%%%%%%%%  scan idx 417  |  folio 406  %%%%%%%%%%

Así, los $P_n$ definidos por las extensiones $A_n$ provienen de un
único torsor $P$ sobre $S$ bajo $H$. Obtenemos una extensión $A$ de
$\Phi$ por $G=A^o$ que algebraiza $\mathcal A$, salvo que resta comprobar
su representabilidad. Consideremos $A$ como un torsor sobre $\Phi$ bajo
$A^o_\Phi$. Este torsor se trivializa después del cambio de base por el
morfismo finito, sobreyectivo y localmente libre
\[
 P\times_S\Phi\longrightarrow\Phi,
\]
pues el cambio de base $P\longrightarrow S$ escinde la extensión $A$
de $\Phi$ por $A^o$. Se puede concluir entonces por uno de los muchos
criterios de Raynaud para la efectividad del descenso
[23, XI~3.1(6)].

\medskip
\noindent
7.2.1.4. Esto termina la demostración de b), y solo queda la afirmación
de plena fidelidad a). Dejamos al lector la tarea de comprobar que esta
afirmación está contenida esencialmente en el método de demostración de
b), o, si lo prefiere, de buscar una demostración directa independiente.

\medskip
\noindent
7.2.1.5. Obsérvese que el recurso al delicado resultado de descenso de
Raynaud resulta innecesario si se supone que $A^o$ es cuasiproyectivo
sobre $S$, o, equivalentemente, que $B$ es proyectivo sobre $S$. Esta
condición se cumple, en particular, cuando se parte del esquema formal
en grupos $\widehat A$ sobre $S$, con $A$ como en 7.1.5. Este es el
único caso en el que utilizaremos 7.2.1.
